\documentclass[10pt,a4paper]{amsart}
\usepackage[margin=19mm]{geometry}
\usepackage{amsmath,amssymb,mathtools}
\usepackage[scr=rsfs,cal=euler]{mathalfa}
\usepackage{xfrac}
\usepackage[unicode,hidelinks,bookmarks=false]{hyperref}
\newcommand{\ie}{i.e.\ }
\newcommand{\cf}{cf.\ }
\newcommand{\resp}{respectively\ }
\newcommand{\Subsection}{\subsection}
\providecommand{\nobreakdash}{\nobreak\hbox{-}}
\setlength{\emergencystretch}{2em}
\multlinegap0pt
\usepackage{bm}
\usepackage{bbm}
\usepackage{dsfont}
\usepackage{xspace}
\usepackage{cancel}
\usepackage{mathtools}
\usepackage{amsthm}
\makeatletter
\@ifundefined{thm}{\newtheorem{thm}{Thm}}{}
\@ifundefined{lem}{\newtheorem{lem}[thm]{Lemma}}{}
\makeatother
\newcommand*\dd{\mathop{}\!\mathrm{d}}
\newcommand{\HH}{\mathbb{H}}
\newcommand{\IS}{\mathbb{S}}
\newcommand{\IT}{\mathbb{T}}
\title[Fully discrete approximation of the semilinear stochastic wa]{Fully discrete approximation of the semilinear stochastic wave equation on the sphere}
\author{David Cohen, Stefano Di Giovacchino, Annika Lang}
\date{Source: arXiv:2602.00556v1 (CC BY 4.0)}
\begin{document}
\maketitle

\noindent\emph{Source and licence.} Fully discrete approximation of the semilinear stochastic wave equation on the sphere. Version arXiv:2602.00556v1 (CC BY 4.0), distributed under the Creative Commons Attribution 4.0 International licence, as recorded in the arXiv licence metadata for this exact version and shown on the abstract page https://arxiv.org/abs/2602.00556v1. Full rights notice: licenses/arxiv-2602.00556-CC-BY-4.0.txt.\par\smallskip

\noindent\textit{Editorial scope.} Counted unit: the Lem whose source label is \texttt{lem\_op2} in the article (source0.tex lines 306--322, source label \texttt{lem\_op2}), delivered together with the article's own complete original proof of it (source0.tex lines 324--369, one \texttt{proof} environment of 46 lines). No mathematics is written, repaired or re-derived: statement and proof are the article's own text line for line. Required context retained uncounted: lem\_op (lines 261--276). Every definition, display and label of those lines is retained unchanged. Omitted from this package and not redistributed: 1--260, 277--305, 323--323, 370--1050. Nothing inside a retained excerpt is omitted or altered: every retained body is delivered exactly as the article's lines are written, so no omission receipt of any kind accompanies this package. Citations: the retained excerpts carry no citation command at all, so no bibliography entry of the article is transcribed and no reference is redistributed. Reference adaptation: none was needed and none was made; every reference inside the delivered unit resolves inside it, so no unresolved reference or citation is produced. Printed slips kept literal: no printed word, symbol, hypothesis or number of the counted statement or of its proof was altered. Layout and class: the article's own class is \texttt{amsart} with packages including inputenc, amssymb, amsmath, graphics, pstricks, tikz, a4wide, rotating, babel, amsthm, bbold, latexsym, color, amsfonts; the wrapper is the lane's pinned \texttt{amsart} editorial wrapper at 10pt on A4, which restates the theorem-like environments the retained text needs and adds the article's own macro definitions that the retained text needs (\texttt{\textbackslash HH}, \texttt{\textbackslash IS}, \texttt{\textbackslash IT}, \texttt{\textbackslash dd}), with extra wrapper packages (bm, bbm, dsfont, xspace, cancel, mathtools). The delivered numbering is the unit's own. Assets: no retained span contains a figure environment, a graphics-inclusion command, a PDF-page inclusion, a table or a TikZ picture; no image file is copied into this package, the package declares no asset dependency and no image file is redistributed with it. Licence: version arXiv:2602.00556v1 of this article is distributed under the Creative Commons Attribution 4.0 International licence, as recorded in the arXiv licence metadata for this exact version and shown on the abstract page https://arxiv.org/abs/2602.00556v1. A full rights notice is delivered at licenses/arxiv-2602.00556-CC-BY-4.0.txt. Characters: every retained excerpt is pure ASCII, so the delivered engine, XeLaTeX with its default Latin Modern text font, reports no missing character.
\begin{lem}
\label{lem_op}
For any $s,t\in\mathbb R$, $u\in H^s(\mathbb{S}^2)$ and $v \in H^{s-1}(\mathbb{S}^2)$, the elements of the group matrix $E(t)$ are bounded by
$$
\begin{aligned}
\left\|C(t)u \right\|_{s} & \le \left\|u\right\|_{s}, \quad 
\left\|\left(-\Delta_{\mathbb{S}^2}\right)^{-\frac{1}{2}}S(t) v \right\|_{s} \le 
\max\left(|t|,\sqrt{\frac32}\right)\left\|v\right\|_{{s-1}}, \quad
\left\|\left(-\Delta_{\mathbb{S}^2}\right)^{\frac{1}{2}} S(t)  u \right\|_{{s-1}} \le \left\|u\right\|_{{s}}.
\end{aligned}
$$
These bounds yield the group estimate 
\begin{equation*} %\label{semi0}
\left\|E(t)\right\|_{\mathcal{L}(\HH^s)}\leq 2\max\left(|t|,\sqrt{\frac32}\right).
\end{equation*}
\end{lem}

\begin{lem}
\label{lem_op2}
For any $s\in\mathbb{R}$, $q \in [0,1]$, $t \in \IT$, $u \in H^{s+q}(\IS^2)$, and $v \in H^{s-1+q}(\IS^2)$, the sine and cosine operators satisfy the bounds
\begin{gather*}
    \| (C(t) - I)u \|_s
         \le 2 t^q \|u\|_{s+q}, \quad
    \| (-\Delta_{\mathbb{S}^2})^{-\frac{1}{2}} S(t)v \|_s
        \le \max\left(T, \sqrt{\tfrac{3}{2}}\right) t^q \|v\|_{s-1+q}, \\
    \left\|(-\Delta_{\mathbb{S}^2})^{\frac{1}{2}}S(t) u\right\|_{s-1}
         \le t^q \left\|u\right\|_{s+q}.
\end{gather*}
These estimates bound the group~$E$ by
\begin{equation*} %\label{semi1}
    \| E(t) - I \|_{\mathcal{L}(\HH^{s+q};\HH^s)}
        \le \max (4, 2T) t^q .
\end{equation*}
\end{lem}

\begin{proof}
    With the same expansion as in the proof of Lemma~\ref{lem_op}, we obtain
    \begin{align*}
        \| (C(t) - I)u \|_s^2
            = \sum_{\ell=0}^\infty \sum_{m=-\ell}^\ell (1+\ell(\ell+1))^{s} (\cos(t(\ell(\ell+1))^{1/2}) - 1)^2 |u^{\ell,m}|^2.
    \end{align*}
    We observe that the summand for $\ell=0$ is zero and that for $\ell \neq 0$ on the one hand 
    \begin{equation*}
        (\cos(t(\ell(\ell+1))^{1/2}) - 1)^2 \le 4
    \end{equation*}
    and on the other hand 
    \begin{equation*}
        (\cos(t(\ell(\ell+1))^{1/2}) - 1)^2
            = \left( \int_0^{t(\ell(\ell+1))^{1/2}} \sin(r) \, \dd r \right)^2
            \le t^2 \ell(\ell+1)
            \le t^2 (1+\ell(\ell+1)),
    \end{equation*}
    since $|\sin(\cdot)|$ is bounded by~$1$. Plugging these bounds into the series expansion therefore yields
    \begin{equation*}
        \| (C(t) - I)u \|_s^2
        \le
        \begin{cases}
            4 \|u\|_s^2 & \text{for } u \in H^s(\IS^2),\\
            t^2 \|u\|_{s+1}^2 & \text{for } u \in H^{s+1}(\IS^2).
        \end{cases}
    \end{equation*}
    The claim follows by interpolation.

    To show the second claim, we use the bound from Lemma~\ref{lem_op} and observe further, with the boundedness of the cosine function, that
    \begin{equation*}
        \sin^2(t(\ell(\ell+1))^{1/2})
            = \left( \int_0^{t(\ell(\ell+1))^{1/2}} \cos(r) \, \dd r\right)^2
            \le t^2 \ell(\ell+1),
    \end{equation*}
    which plugged into the series expansion in the proof of Lemma~\ref{lem_op} yields
    \begin{equation*}
        \| (-\Delta_{\mathbb{S}^2})^{-\frac{1}{2}} S(t)v \|_s^2
        \le
        \begin{cases}
            \max\left(t^2,\tfrac{3}{2}\right) \|v\|_{s-1}^2 & \text{for } v \in H^{s-1}(\IS^2),\\
            t^2 \|v\|_{s}^2 & \text{for } v \in H^{s}(\IS^2).
        \end{cases}
    \end{equation*}
    The claim follows by interpolation and by bounding $t$ by~$T$.
    The last claim follows analogously, and its proof is therefore omitted.
\end{proof}

\end{document}
